\clearpage
\VerdierPhysicalPageAnchor{0152}%
\begingroup
\small
\setlength{\parskip}{0pt}%
\setlength{\emergencystretch}{1.5em}%

\VerdierSemanticAnchor{ch3}
\section{三角范畴的阿贝尔化。}

\VerdierSemanticAnchor{ch3-3-1}
\subsection{范畴 \texorpdfstring{$\mathrm A(\mathcal D)$}{A(D)} 的构造。}

\VerdierSemanticAnchor{ch3-3-1-1}
\noindent\textbf{3.1.1.}\quad
设 $\mathcal D$ 为加法范畴，$\mathcal{F}\ell_1(\mathcal D)$ 为
$\mathcal D$ 的态射范畴（第一章，3.3.6）。考虑下列交换图：

\VerdierSemanticAnchor{ch3-3-1-1-square}
\[
\begin{tikzcd}[column sep=4.2em,row sep=2.5em]
X \arrow[r,"f"] \arrow[d,"g"'] & Y \arrow[d,"h"] \\
X' \arrow[r,"f'"'] & Y'
\end{tikzcd}
\]

\VerdierSemanticAnchor{ch3-3-1-1-negligible-morphism}
\noindent 上述 $\mathcal{F}\ell_1(\mathcal D)$ 中的态射在
$f'g=hf=0$ 成立时称为\emph{可忽略态射}。可忽略态射与任意态射
按任一顺序复合，所得态射仍可忽略；两个可忽略态射的直和仍可忽略。

\VerdierSemanticAnchor{ch3-3-1-1-quotient-category}
以 $\mathrm A(\mathcal D)$ 表示如下范畴：其对象就是
$\mathcal{F}\ell_1(\mathcal D)$ 的对象，其态射则是
$\mathcal{F}\ell_1(\mathcal D)$ 的态射模去可忽略态射。换言之，对
$\mathrm A(\mathcal D)$ 的任意对象 $U,V$，群
$\operatorname{Hom}_{\mathrm A(\mathcal D)}(U,V)$ 是
$\operatorname{Hom}_{\mathcal{F}\ell_1(\mathcal D)}(U,V)$ 对可忽略态射
子群的商。范畴 $\mathrm A(\mathcal D)$ 是加法范畴。

\VerdierSemanticAnchor{ch3-3-1-1-canonical-quotient}
设 $C$ 为一个范畴。与典范函子
\[
\mathcal{F}\ell_1(\mathcal D)\longrightarrow\mathrm A(\mathcal D)
\]
复合给出函子

\VerdierSemanticAnchor{ch3-3-1-1-1}
\[
\mathcal H\!om(\mathrm A(\mathcal D),C)
\longrightarrow
\mathcal H\!om(\mathcal{F}\ell_1(\mathcal D),C).
\tag{3.1.1.1}
\]
该函子充分忠实，并且在对象上单射。它的像是如下满子范畴：其对象为从
$\mathcal{F}\ell_1(\mathcal D)$ 到 $C$ 的函子，这些函子把相差一个可忽略态射的
态射送到相等的态射。

\VerdierSemanticAnchor{ch3-3-1-2}
\medskip
\noindent\textbf{3.1.2.}\quad
从 $\mathcal D$ 到 $\mathcal{F}\ell_1(\mathcal D)$ 的典范函子
$X\longmapsto(X\xrightarrow{\operatorname{id}_X}X)$ 与典范函子
$\mathcal{F}\ell_1(\mathcal D)\longrightarrow\mathrm A(\mathcal D)$ 复合，所得函子记为

\VerdierSemanticAnchor{ch3-3-1-2-1}
\[
\mathcal H\colon\mathcal D\longrightarrow\mathrm A(\mathcal D).
\tag{3.1.2.1}
\]
函子 $\mathcal H$ 是加法函子；它充分忠实，并且在对象上单射。因此，它把
$\mathcal D$ 实现为 $\mathrm A(\mathcal D)$ 的满子范畴。
\endgroup
